\section{Conclusion}
This proposal outlines a neuro-symbolic decoding framework that integrates a
compiled Syntax Oracle as an abstract machine into the generation loop of a language 
model, supported by Oracle-guided fine-tuning. The objective is to enforce 
structural and semantic correctness without collapsing generation
flexibility or incurring prohibitive inference-time overhead. By combining
token-syntax-level masking, element-level semantic validation, and targeted LoRA
adaptation, the system aims to shift part of the verification burden from the
symbolic engine into the model's predictive distribution.

Naturally, several difficulties are expected. The interaction between the Oracle and the
model's internal priors may produce unstable decoding behaviour, especially in
schemas with deep dependency chains. Also, rollback mechanisms, while necessary, risk
introducing latency spikes if the model repeatedly proposes invalid spans. The
quality and coverage of Oracle-generated training data will directly affect how
well the model internalizes constraints, and poorly balanced datasets may lead
to brittle or overfitted adapters. Finally, co-ordinating a C++ verification
engine with Python-based inference introduces practical engineering challenges
around synchronization and memory access that may limit throughput.

Even with these limitations, the approach provides a concrete path toward
quantifying and reducing the symbolic workload during structured generation.
Rather than treating verification as an external filter, the framework treats it
as an active component of the decoding process and a source of training signal.

If successful, this work will demonstrate that formally constrained generation
can be achieved through a tighter coupling between neural prediction and
symbolic structure, offering a more reliable alternative to prompt-based
heuristics and post-hoc validation.